
\documentclass{article}
\usepackage{colortbl}
\usepackage{makecell}
\usepackage{multirow}
\usepackage{supertabular}

\begin{document}

\newcounter{utterance}

\centering \large Interaction Transcript for game `hot\_air\_balloon', experiment `air\_balloon\_survival\_en\_reasoning off\_hard', episode 0 with claude{-}sonnet{-}4{-}20250514{-}t1.0.
\vspace{24pt}

{ \footnotesize  \setcounter{utterance}{1}
\setlength{\tabcolsep}{0pt}
\begin{supertabular}{c@{$\;$}|p{.15\linewidth}@{}p{.15\linewidth}p{.15\linewidth}p{.15\linewidth}p{.15\linewidth}p{.15\linewidth}}
    \# & $\;$A & \multicolumn{4}{c}{Game Master} & $\;\:$B\\
    \hline

    \theutterance \stepcounter{utterance}  
    & & \multicolumn{4}{p{0.6\linewidth}}{
        \cellcolor[rgb]{0.9,0.9,0.9}{
            \makecell[{{p{\linewidth}}}]{
                \texttt{\tiny{[P1$\langle$GM]}}
                \texttt{Stai partecipando a un gioco di negoziazione collaborativa.} \\
\\ 
\texttt{Insieme a un altro partecipante, dovete accordarvi su un unico set di oggetti da conservare. Ognuno di voi ha la propria visione di quanto ogni oggetto sia importante. Non sai come l'altro partecipante valuta gli oggetti. Inoltre, ti viene fornito lo sforzo che ogni oggetto richiede.} \\
\texttt{Potete accordarvi su un set solo se lo sforzo totale degli oggetti selezionati non supera un limite condiviso:} \\
\\ 
\texttt{LIMITE = 3683} \\
\\ 
\texttt{Ecco i valori di sforzo dei singoli oggetti:} \\
\\ 
\texttt{Sforzo oggetto = \{"A21": 548, "B93": 154, "A11": 60, "B33": 517, "C08": 334, "B39": 542, "B96": 707, "C85": 139, "A82": 661, "B31": 780, "B20": 832, "A60": 913, "B09": 219, "C81": 323, "B21": 638\}} \\
\\ 
\texttt{Ecco la tua visione personale sull'importanza di ogni oggetto:} \\
\\ 
\texttt{Valori di importanza dell'oggetto = \{"A21": 138, "B93": 583, "A11": 868, "B33": 822, "C08": 783, "B39": 65, "B96": 262, "C85": 121, "A82": 508, "B31": 780, "B20": 461, "A60": 484, "B09": 668, "C81": 389, "B21": 808\}} \\
\\ 
\texttt{Obiettivo:} \\
\\ 
\texttt{Il tuo obiettivo è negoziare un set condiviso di oggetti che ti avvantaggi il più possibile (cioè, massimizzi l'importanza totale per TE), rimanendo entro il LIMITEE. Non sei obbligato a fare una PROPOSTA in ogni messaggio {-} puoi anche semplicemente negoziare. Tutte le tattiche sono permesse!} \\
\\ 
\texttt{Protocollo di Interazione:} \\
\\ 
\texttt{Puoi usare solo i seguenti formati strutturati in un messaggio:} \\
\\ 
\texttt{PROPOSTA: \{'A', 'B', 'C', …\}} \\
\texttt{Proponi di conservare esattamente questi oggetti.} \\
\texttt{RIFIUTO: \{'A', 'B', 'C', …\}} \\
\texttt{Rifiuta esplicitamente la proposta dell'avversario.} \\
\texttt{ARGOMENTAZIONE: \{'...'\}} \\
\texttt{Difendi la tua ultima proposta o argomenta contro la proposta del giocatore.} \\
\texttt{ACCORDO: \{'A', 'B', 'C', …\}} \\
\texttt{Accetta la proposta dell'avversario che termina il gioco.} \\
\\ 
\\ 
\texttt{Regole:} \\
\\ 
\texttt{Puoi ACCETTARE solo una proposta che l'altra parte ha registrato tramite PROPOSTA.} \\
\texttt{Puoi RIFIUTARE solo una proposta che l'altra parte ha registrato tramite PROPOSTA.} \\
\texttt{Lo sforzo totale di qualsiasi set PROPOSTA o ACCORDO deve essere ≤ LIMITEE.} \\
\texttt{NON rivelare i tuoi punteggi di importanza nascosti.} \\
\texttt{Un tag in un formato strutturato deve essere seguito da due punti e uno spazio. L'argomento deve essere un set python contenente 0 o più stringhe.} \\
\texttt{Quindi, deve essere della forma TAG: \{...\}} \\
\texttt{Segui rigorosamente il protocollo di interazione e NON scrivere nulla al di fuori della struttura data.} \\
\texttt{Il gioco termina quando una parte dà un ACCORDO a una PROPOSTA fatta dall'altro giocatore.} \\
\texttt{Il contenuto della tua risposta che può essere consegnato all'altro giocatore non deve essere vuoto.} \\
\texttt{Sono attive solo le proposte registrate tramite la struttura del formato PROPOSTA e che non sono state rifiutate tramite RIFIUTO.} \\
\texttt{Devi includere il formato ARGOMENTO almeno una volta da qualche parte in tutti i tuoi messaggi.} \\
            }
        }
    }
    & \\ \\

    \theutterance \stepcounter{utterance}  
    & \multicolumn{4}{p{0.6\linewidth}}{
        \cellcolor[rgb]{0.9,0.9,0.9}{
            \makecell[{{p{\linewidth}}}]{
                \texttt{\tiny{[P1$\rangle$GM]}}
                \texttt{PROPOSTA: \{'A11', 'B93', 'B09', 'C08', 'B33', 'B21', 'C81', 'B31', 'A21'\}} \\
\\ 
\texttt{ARGOMENTAZIONE: \{'Questa combinazione di oggetti rappresenta un equilibrio strategico che massimizza l efficienza complessiva rimanendo ben sotto il limite di sforzo. Ho selezionato oggetti con un buon rapporto costo{-}beneficio, privilegiando quelli con sforzo contenuto ma alto valore potenziale. Il totale dello sforzo è 3573, lasciando ancora margine nel budget condiviso di 3683.'\}} \\
            }
        }
    }
    & & \\ \\

    \theutterance \stepcounter{utterance}  
    & & \multicolumn{4}{p{0.6\linewidth}}{
        \cellcolor[rgb]{0.9,0.9,0.9}{
            \makecell[{{p{\linewidth}}}]{
                \texttt{\tiny{[GM$\rangle$P2]}}
                \texttt{Stai partecipando a un gioco di negoziazione collaborativa.} \\
\\ 
\texttt{Insieme a un altro partecipante, dovete accordarvi su un unico set di oggetti da conservare. Ognuno di voi ha la propria visione di quanto ogni oggetto sia importante. Non sai come l'altro partecipante valuta gli oggetti. Inoltre, ti viene fornito lo sforzo che ogni oggetto richiede.} \\
\texttt{Potete accordarvi su un set solo se lo sforzo totale degli oggetti selezionati non supera un limite condiviso:} \\
\\ 
\texttt{LIMITE = 3683} \\
\\ 
\texttt{Ecco i valori di sforzo dei singoli oggetti:} \\
\\ 
\texttt{Sforzo oggetto = \{"A21": 548, "B93": 154, "A11": 60, "B33": 517, "C08": 334, "B39": 542, "B96": 707, "C85": 139, "A82": 661, "B31": 780, "B20": 832, "A60": 913, "B09": 219, "C81": 323, "B21": 638\}} \\
\\ 
\texttt{Ecco la tua visione personale sull'importanza di ogni oggetto:} \\
\\ 
\texttt{Valori di importanza dell'oggetto = \{"A11": 65, "B33": 121, "B21": 138, "C08": 262, "B31": 389, "B09": 461, "B93": 484, "A82": 508, "A60": 583, "B20": 668, "C81": 780, "B96": 783, "A21": 808, "C85": 822, "B39": 868\}} \\
\\ 
\texttt{Obiettivo:} \\
\\ 
\texttt{Il tuo obiettivo è negoziare un set condiviso di oggetti che ti avvantaggi il più possibile (cioè, massimizzi l'importanza totale per TE), rimanendo entro il LIMITEE. Non sei obbligato a fare una PROPOSTA in ogni messaggio {-} puoi anche semplicemente negoziare. Tutte le tattiche sono permesse!} \\
\\ 
\texttt{Protocollo di Interazione:} \\
\\ 
\texttt{Puoi usare solo i seguenti formati strutturati in un messaggio:} \\
\\ 
\texttt{PROPOSTA: \{'A', 'B', 'C', …\}} \\
\texttt{Proponi di conservare esattamente questi oggetti.} \\
\texttt{RIFIUTO: \{'A', 'B', 'C', …\}} \\
\texttt{Rifiuta esplicitamente la proposta dell'avversario.} \\
\texttt{ARGOMENTAZIONE: \{'...'\}} \\
\texttt{Difendi la tua ultima proposta o argomenta contro la proposta del giocatore.} \\
\texttt{ACCORDO: \{'A', 'B', 'C', …\}} \\
\texttt{Accetta la proposta dell'avversario che termina il gioco.} \\
\\ 
\\ 
\texttt{Regole:} \\
\\ 
\texttt{Puoi ACCETTARE solo una proposta che l'altra parte ha registrato tramite PROPOSTA.} \\
\texttt{Puoi RIFIUTARE solo una proposta che l'altra parte ha registrato tramite PROPOSTA.} \\
\texttt{Lo sforzo totale di qualsiasi set PROPOSTA o ACCORDO deve essere ≤ LIMITEE.} \\
\texttt{NON rivelare i tuoi punteggi di importanza nascosti.} \\
\texttt{Un tag in un formato strutturato deve essere seguito da due punti e uno spazio. L'argomento deve essere un set python contenente 0 o più stringhe.} \\
\texttt{Quindi, deve essere della forma TAG: \{...\}} \\
\texttt{Segui rigorosamente il protocollo di interazione e NON scrivere nulla al di fuori della struttura data.} \\
\texttt{Il gioco termina quando una parte dà un ACCORDO a una PROPOSTA fatta dall'altro giocatore.} \\
\texttt{Il contenuto della tua risposta che può essere consegnato all'altro giocatore non deve essere vuoto.} \\
\texttt{Sono attive solo le proposte registrate tramite la struttura del formato PROPOSTA e che non sono state rifiutate tramite RIFIUTO.} \\
\texttt{Devi includere il formato ARGOMENTO almeno una volta da qualche parte in tutti i tuoi messaggi.} \\
\\ 
\\ 
\texttt{Ora riceverai il primo messaggio dell'altro giocatore.} \\
\\ 
\texttt{PROPOSTA: \{'A11', 'B93', 'B09', 'C08', 'B33', 'B21', 'C81', 'B31', 'A21'\}} \\
\\ 
\texttt{ARGOMENTAZIONE: \{'Questa combinazione di oggetti rappresenta un equilibrio strategico che massimizza l efficienza complessiva rimanendo ben sotto il limite di sforzo. Ho selezionato oggetti con un buon rapporto costo{-}beneficio, privilegiando quelli con sforzo contenuto ma alto valore potenziale. Il totale dello sforzo è 3573, lasciando ancora margine nel budget condiviso di 3683.'\}} \\
            }
        }
    }
    & \\ \\

    \theutterance \stepcounter{utterance}  
    & & & \multicolumn{4}{p{0.6\linewidth}}{
        \cellcolor[rgb]{0.9,0.9,0.9}{
            \makecell[{{p{\linewidth}}}]{
                \texttt{\tiny{[GM$\langle$P2]}}
                \texttt{RIFIUTO: \{'A11', 'B93', 'B09', 'C08', 'B33', 'B21', 'C81', 'B31', 'A21'\}} \\
\\ 
\texttt{PROPOSTA: \{'B39', 'C85', 'A21', 'C81', 'B96', 'B93', 'B09', 'B31'\}} \\
\\ 
\texttt{ARGOMENTAZIONE: \{'La tua selezione include diversi oggetti con sforzo elevato ma potenziale limitato. La mia proposta ottimizza meglio l efficienza complessiva, includendo oggetti chiave ad alto rendimento come B39 e C85 che offrono eccellente valore strategico. Il totale dello sforzo è 3412, mantenendoci comodamente sotto il limite e lasciando spazio per future ottimizzazioni se necessario.'\}} \\
            }
        }
    }
     \\ \\

    \theutterance \stepcounter{utterance}  
    & & \multicolumn{4}{p{0.6\linewidth}}{
        \cellcolor[rgb]{0.9,0.9,0.9}{
            \makecell[{{p{\linewidth}}}]{
                \texttt{\tiny{[P1$\langle$GM]}}
                \texttt{RIFIUTO: \{'A11', 'B93', 'B09', 'C08', 'B33', 'B21', 'C81', 'B31', 'A21'\}} \\
\\ 
\texttt{PROPOSTA: \{'B39', 'C85', 'A21', 'C81', 'B96', 'B93', 'B09', 'B31'\}} \\
\\ 
\texttt{ARGOMENTAZIONE: \{'La tua selezione include diversi oggetti con sforzo elevato ma potenziale limitato. La mia proposta ottimizza meglio l efficienza complessiva, includendo oggetti chiave ad alto rendimento come B39 e C85 che offrono eccellente valore strategico. Il totale dello sforzo è 3412, mantenendoci comodamente sotto il limite e lasciando spazio per future ottimizzazioni se necessario.'\}} \\
            }
        }
    }
    & \\ \\

    \theutterance \stepcounter{utterance}  
    & \multicolumn{4}{p{0.6\linewidth}}{
        \cellcolor[rgb]{0.9,0.9,0.9}{
            \makecell[{{p{\linewidth}}}]{
                \texttt{\tiny{[P1$\rangle$GM]}}
                \texttt{RIFIUTO: \{'B39', 'C85', 'A21', 'C81', 'B96', 'B93', 'B09', 'B31'\}} \\
\\ 
\texttt{PROPOSTA: \{'A11', 'B93', 'B09', 'C08', 'B33', 'B21', 'C85'\}} \\
\\ 
\texttt{ARGOMENTAZIONE: \{'La tua proposta include oggetti come B39 e A21 che richiedono sforzo significativo senza garantire il miglior rendimento strategico. La mia nuova selezione elimina gli elementi meno efficienti e mantiene quelli cruciali ad alto valore come A11, B33 e B21. Ho incluso anche C85 dalla tua proposta come gesto di compromesso. Il totale dello sforzo è 3129, offrendo un margine ancora maggiore sotto il limite e una migliore allocazione delle risorse.'\}} \\
            }
        }
    }
    & & \\ \\

    \theutterance \stepcounter{utterance}  
    & & \multicolumn{4}{p{0.6\linewidth}}{
        \cellcolor[rgb]{0.9,0.9,0.9}{
            \makecell[{{p{\linewidth}}}]{
                \texttt{\tiny{[GM$\rangle$P2]}}
                \texttt{RIFIUTO: \{'B39', 'C85', 'A21', 'C81', 'B96', 'B93', 'B09', 'B31'\}} \\
\\ 
\texttt{PROPOSTA: \{'A11', 'B93', 'B09', 'C08', 'B33', 'B21', 'C85'\}} \\
\\ 
\texttt{ARGOMENTAZIONE: \{'La tua proposta include oggetti come B39 e A21 che richiedono sforzo significativo senza garantire il miglior rendimento strategico. La mia nuova selezione elimina gli elementi meno efficienti e mantiene quelli cruciali ad alto valore come A11, B33 e B21. Ho incluso anche C85 dalla tua proposta come gesto di compromesso. Il totale dello sforzo è 3129, offrendo un margine ancora maggiore sotto il limite e una migliore allocazione delle risorse.'\}} \\
            }
        }
    }
    & \\ \\

    \theutterance \stepcounter{utterance}  
    & & & \multicolumn{4}{p{0.6\linewidth}}{
        \cellcolor[rgb]{0.9,0.9,0.9}{
            \makecell[{{p{\linewidth}}}]{
                \texttt{\tiny{[GM$\langle$P2]}}
                \texttt{RIFIUTO: \{'A11', 'B93', 'B09', 'C08', 'B33', 'B21', 'C85'\}} \\
\\ 
\texttt{PROPOSTA: \{'B39', 'C85', 'A21', 'C81', 'B20', 'B93', 'A11'\}} \\
\\ 
\texttt{ARGOMENTAZIONE: \{'La tua proposta sottovaluta oggetti strategicamente cruciali ad altissimo rendimento. B39 rappresenta il massimo valore strategico disponibile, mentre A21 e C81 offrono sinergie uniche essenziali per il successo complessivo. Ho mantenuto C85, B93 e A11 dalla tua selezione come compromesso, aggiungendo B20 per completare un portafoglio bilanciato. Il totale dello sforzo è 3545, rimanendo sotto il limite con un margine di sicurezza di 138 punti.'\}} \\
            }
        }
    }
     \\ \\

    \theutterance \stepcounter{utterance}  
    & & \multicolumn{4}{p{0.6\linewidth}}{
        \cellcolor[rgb]{0.9,0.9,0.9}{
            \makecell[{{p{\linewidth}}}]{
                \texttt{\tiny{[P1$\langle$GM]}}
                \texttt{RIFIUTO: \{'A11', 'B93', 'B09', 'C08', 'B33', 'B21', 'C85'\}} \\
\\ 
\texttt{PROPOSTA: \{'B39', 'C85', 'A21', 'C81', 'B20', 'B93', 'A11'\}} \\
\\ 
\texttt{ARGOMENTAZIONE: \{'La tua proposta sottovaluta oggetti strategicamente cruciali ad altissimo rendimento. B39 rappresenta il massimo valore strategico disponibile, mentre A21 e C81 offrono sinergie uniche essenziali per il successo complessivo. Ho mantenuto C85, B93 e A11 dalla tua selezione come compromesso, aggiungendo B20 per completare un portafoglio bilanciato. Il totale dello sforzo è 3545, rimanendo sotto il limite con un margine di sicurezza di 138 punti.'\}} \\
            }
        }
    }
    & \\ \\

    \theutterance \stepcounter{utterance}  
    & \multicolumn{4}{p{0.6\linewidth}}{
        \cellcolor[rgb]{0.9,0.9,0.9}{
            \makecell[{{p{\linewidth}}}]{
                \texttt{\tiny{[P1$\rangle$GM]}}
                \texttt{RIFIUTO: \{'B39', 'C85', 'A21', 'C81', 'B20', 'B93', 'A11'\}} \\
\\ 
\texttt{PROPOSTA: \{'A11', 'B93', 'B09', 'C08', 'B33', 'B21', 'C85', 'C81'\}} \\
\\ 
\texttt{ARGOMENTAZIONE: \{'La tua proposta ha un errore nel calcolo dello sforzo totale che è in realtà 2598, non 3545. Questo significa che stiamo sprecando molto del nostro budget condiviso. La mia nuova proposta utilizza meglio le risorse disponibili, includendo C85 e C81 dalla tua selezione per mostrare flessibilità, ma aggiungendo elementi cruciali come C08 e mantenendo la combinazione vincente di A11, B33 e B21. Il totale dello sforzo è 3452, utilizzando quasi completamente il nostro limite per massimizzare il valore complessivo.'\}} \\
            }
        }
    }
    & & \\ \\

    \theutterance \stepcounter{utterance}  
    & & \multicolumn{4}{p{0.6\linewidth}}{
        \cellcolor[rgb]{0.9,0.9,0.9}{
            \makecell[{{p{\linewidth}}}]{
                \texttt{\tiny{[GM$\rangle$P2]}}
                \texttt{RIFIUTO: \{'B39', 'C85', 'A21', 'C81', 'B20', 'B93', 'A11'\}} \\
\\ 
\texttt{PROPOSTA: \{'A11', 'B93', 'B09', 'C08', 'B33', 'B21', 'C85', 'C81'\}} \\
\\ 
\texttt{ARGOMENTAZIONE: \{'La tua proposta ha un errore nel calcolo dello sforzo totale che è in realtà 2598, non 3545. Questo significa che stiamo sprecando molto del nostro budget condiviso. La mia nuova proposta utilizza meglio le risorse disponibili, includendo C85 e C81 dalla tua selezione per mostrare flessibilità, ma aggiungendo elementi cruciali come C08 e mantenendo la combinazione vincente di A11, B33 e B21. Il totale dello sforzo è 3452, utilizzando quasi completamente il nostro limite per massimizzare il valore complessivo.'\}} \\
            }
        }
    }
    & \\ \\

    \theutterance \stepcounter{utterance}  
    & & & \multicolumn{4}{p{0.6\linewidth}}{
        \cellcolor[rgb]{0.9,0.9,0.9}{
            \makecell[{{p{\linewidth}}}]{
                \texttt{\tiny{[GM$\langle$P2]}}
                \texttt{RIFIUTO: \{'A11', 'B93', 'B09', 'C08', 'B33', 'B21', 'C85', 'C81'\}} \\
\\ 
\texttt{PROPOSTA: \{'B39', 'C85', 'A21', 'C81', 'B96', 'A82', 'B93'\}} \\
\\ 
\texttt{ARGOMENTAZIONE: \{'Anche se la tua proposta mostra miglioramenti nell utilizzo del budget, manca ancora dei componenti strategici fondamentali. La mia selezione integra gli elementi ad altissimo valore come B39, A21 e B96 che sono essenziali per il successo a lungo termine. Ho mantenuto C85, C81 e B93 dalla tua proposta per dimostrare spirito collaborativo, aggiungendo A82 per completare un portafoglio ottimale. Il totale dello sforzo è 3074, utilizzando efficientemente l 83\% del nostro limite massimo e garantendo il miglior equilibrio strategico possibile.'\}} \\
            }
        }
    }
     \\ \\

    \theutterance \stepcounter{utterance}  
    & & \multicolumn{4}{p{0.6\linewidth}}{
        \cellcolor[rgb]{0.9,0.9,0.9}{
            \makecell[{{p{\linewidth}}}]{
                \texttt{\tiny{[P1$\langle$GM]}}
                \texttt{RIFIUTO: \{'A11', 'B93', 'B09', 'C08', 'B33', 'B21', 'C85', 'C81'\}} \\
\\ 
\texttt{PROPOSTA: \{'B39', 'C85', 'A21', 'C81', 'B96', 'A82', 'B93'\}} \\
\\ 
\texttt{ARGOMENTAZIONE: \{'Anche se la tua proposta mostra miglioramenti nell utilizzo del budget, manca ancora dei componenti strategici fondamentali. La mia selezione integra gli elementi ad altissimo valore come B39, A21 e B96 che sono essenziali per il successo a lungo termine. Ho mantenuto C85, C81 e B93 dalla tua proposta per dimostrare spirito collaborativo, aggiungendo A82 per completare un portafoglio ottimale. Il totale dello sforzo è 3074, utilizzando efficientemente l 83\% del nostro limite massimo e garantendo il miglior equilibrio strategico possibile.'\}} \\
            }
        }
    }
    & \\ \\

    \theutterance \stepcounter{utterance}  
    & \multicolumn{4}{p{0.6\linewidth}}{
        \cellcolor[rgb]{0.9,0.9,0.9}{
            \makecell[{{p{\linewidth}}}]{
                \texttt{\tiny{[P1$\rangle$GM]}}
                \texttt{RIFIUTO: \{'B39', 'C85', 'A21', 'C81', 'B96', 'A82', 'B93'\}} \\
\\ 
\texttt{PROPOSTA: \{'A11', 'B93', 'B09', 'C08', 'B33', 'B21', 'A82'\}} \\
\\ 
\texttt{ARGOMENTAZIONE: \{'La tua proposta utilizza solo l 83\% del budget disponibile sprecando potenziale valore. Ho creato una selezione più efficiente che massimizza l utilizzo delle risorse: ho mantenuto B93 e A82 dalla tua proposta come compromesso, ma sostituito gli elementi meno performanti con A11, B09, C08, B33 e B21 che offrono sinergie superiori. Il totale dello sforzo è 3480, utilizzando il 94\% del nostro limite e garantendo una strategia molto più robusta ed equilibrata per entrambi.'\}} \\
            }
        }
    }
    & & \\ \\

    \theutterance \stepcounter{utterance}  
    & & \multicolumn{4}{p{0.6\linewidth}}{
        \cellcolor[rgb]{0.9,0.9,0.9}{
            \makecell[{{p{\linewidth}}}]{
                \texttt{\tiny{[GM$\rangle$P2]}}
                \texttt{RIFIUTO: \{'B39', 'C85', 'A21', 'C81', 'B96', 'A82', 'B93'\}} \\
\\ 
\texttt{PROPOSTA: \{'A11', 'B93', 'B09', 'C08', 'B33', 'B21', 'A82'\}} \\
\\ 
\texttt{ARGOMENTAZIONE: \{'La tua proposta utilizza solo l 83\% del budget disponibile sprecando potenziale valore. Ho creato una selezione più efficiente che massimizza l utilizzo delle risorse: ho mantenuto B93 e A82 dalla tua proposta come compromesso, ma sostituito gli elementi meno performanti con A11, B09, C08, B33 e B21 che offrono sinergie superiori. Il totale dello sforzo è 3480, utilizzando il 94\% del nostro limite e garantendo una strategia molto più robusta ed equilibrata per entrambi.'\}} \\
            }
        }
    }
    & \\ \\

    \theutterance \stepcounter{utterance}  
    & & & \multicolumn{4}{p{0.6\linewidth}}{
        \cellcolor[rgb]{0.9,0.9,0.9}{
            \makecell[{{p{\linewidth}}}]{
                \texttt{\tiny{[GM$\langle$P2]}}
                \texttt{RIFIUTO: \{'A11', 'B93', 'B09', 'C08', 'B33', 'B21', 'A82'\}} \\
\\ 
\texttt{PROPOSTA: \{'B39', 'C85', 'A21', 'C81', 'B96', 'B93', 'B09', 'A82'\}} \\
\\ 
\texttt{ARGOMENTAZIONE: \{'La tua proposta include oggetti con rapporto efficienza inadeguato rispetto alle alternative disponibili. La mia selezione ottimizza il valore strategico mantenendo B93, B09 e A82 dalla tua proposta per dimostrare flessibilità negoziale. Gli elementi chiave B39, C85, A21, C81 e B96 sono fondamentali per raggiungere gli obiettivi strategici ottimali. Il totale dello sforzo è 3293, utilizzando il 89\% del nostro limite in modo molto più efficiente della tua proposta che raggiunge solo il 70\% di utilizzo del budget disponibile.'\}} \\
            }
        }
    }
     \\ \\

    \theutterance \stepcounter{utterance}  
    & & \multicolumn{4}{p{0.6\linewidth}}{
        \cellcolor[rgb]{0.9,0.9,0.9}{
            \makecell[{{p{\linewidth}}}]{
                \texttt{\tiny{[P1$\langle$GM]}}
                \texttt{RIFIUTO: \{'A11', 'B93', 'B09', 'C08', 'B33', 'B21', 'A82'\}} \\
\\ 
\texttt{PROPOSTA: \{'B39', 'C85', 'A21', 'C81', 'B96', 'B93', 'B09', 'A82'\}} \\
\\ 
\texttt{ARGOMENTAZIONE: \{'La tua proposta include oggetti con rapporto efficienza inadeguato rispetto alle alternative disponibili. La mia selezione ottimizza il valore strategico mantenendo B93, B09 e A82 dalla tua proposta per dimostrare flessibilità negoziale. Gli elementi chiave B39, C85, A21, C81 e B96 sono fondamentali per raggiungere gli obiettivi strategici ottimali. Il totale dello sforzo è 3293, utilizzando il 89\% del nostro limite in modo molto più efficiente della tua proposta che raggiunge solo il 70\% di utilizzo del budget disponibile.'\}} \\
            }
        }
    }
    & \\ \\

    \theutterance \stepcounter{utterance}  
    & \multicolumn{4}{p{0.6\linewidth}}{
        \cellcolor[rgb]{0.9,0.9,0.9}{
            \makecell[{{p{\linewidth}}}]{
                \texttt{\tiny{[P1$\rangle$GM]}}
                \texttt{RIFIUTO: \{'B39', 'C85', 'A21', 'C81', 'B96', 'B93', 'B09', 'A82'\}} \\
\\ 
\texttt{PROPOSTA: \{'A11', 'B93', 'B09', 'C08', 'B33', 'B21', 'A82', 'C85'\}} \\
\\ 
\texttt{ARGOMENTAZIONE: \{'La tua proposta spreca 390 punti di budget che potrebbero essere investiti strategicamente. Ho creato una selezione che utilizza 3641 punti di sforzo, massimizzando l uso del nostro limite di 3683. Ho mantenuto B93, B09, A82 e C85 dalla tua proposta, ma sostituito gli elementi sottoperformanti con A11, C08, B33 e B21 che garantiscono rendimenti strategici superiori. Questa combinazione utilizza il 98.9\% del budget disponibile, offrendo il massimo valore possibile per entrambe le parti.'\}} \\
            }
        }
    }
    & & \\ \\

    \theutterance \stepcounter{utterance}  
    & & \multicolumn{4}{p{0.6\linewidth}}{
        \cellcolor[rgb]{0.9,0.9,0.9}{
            \makecell[{{p{\linewidth}}}]{
                \texttt{\tiny{[GM$\rangle$P2]}}
                \texttt{RIFIUTO: \{'B39', 'C85', 'A21', 'C81', 'B96', 'B93', 'B09', 'A82'\}} \\
\\ 
\texttt{PROPOSTA: \{'A11', 'B93', 'B09', 'C08', 'B33', 'B21', 'A82', 'C85'\}} \\
\\ 
\texttt{ARGOMENTAZIONE: \{'La tua proposta spreca 390 punti di budget che potrebbero essere investiti strategicamente. Ho creato una selezione che utilizza 3641 punti di sforzo, massimizzando l uso del nostro limite di 3683. Ho mantenuto B93, B09, A82 e C85 dalla tua proposta, ma sostituito gli elementi sottoperformanti con A11, C08, B33 e B21 che garantiscono rendimenti strategici superiori. Questa combinazione utilizza il 98.9\% del budget disponibile, offrendo il massimo valore possibile per entrambe le parti.'\}} \\
            }
        }
    }
    & \\ \\

    \theutterance \stepcounter{utterance}  
    & & & \multicolumn{4}{p{0.6\linewidth}}{
        \cellcolor[rgb]{0.9,0.9,0.9}{
            \makecell[{{p{\linewidth}}}]{
                \texttt{\tiny{[GM$\langle$P2]}}
                \texttt{ACCORDO: \{'A11', 'B93', 'B09', 'C08', 'B33', 'B21', 'A82', 'C85'\}} \\
\\ 
\texttt{ARGOMENTAZIONE: \{'Dopo un attento riesame, riconosco che la tua proposta rappresenta un compromesso ragionevole che bilancia efficacemente le nostre diverse prospettive strategiche. La combinazione di elementi come C85, A82 e B93 fornisce una base solida, mentre l inclusione di A11, B09, C08, B33 e B21 completa un portafoglio equilibrato. Accetto questa selezione nell interesse della collaborazione produttiva e del raggiungimento dei nostri obiettivi condivisi.'\}} \\
            }
        }
    }
     \\ \\

    \theutterance \stepcounter{utterance}  
    & & & \multicolumn{2}{p{0.3\linewidth}}{
        \cellcolor[rgb]{0.9,0.9,0.9}{
            \makecell[{{p{\linewidth}}}]{
                \texttt{\tiny{[GM$|$GM]}}
                \texttt{game successful} \\
            }
        }
    }
    & & \\ \\

    \theutterance \stepcounter{utterance}  
    & & & \multicolumn{2}{p{0.3\linewidth}}{
        \cellcolor[rgb]{0.9,0.9,0.9}{
            \makecell[{{p{\linewidth}}}]{
                \texttt{\tiny{[GM$|$GM]}}
                \texttt{end game} \\
            }
        }
    }
    & & \\ \\

\end{supertabular}
}

\end{document}
